Fix an infinite set $A\subset\mathbb{R}^2$. Assume:

$\exists\,\epsilon>0$ such that $\forall n\in\mathbb{N}_{\ge 1}$ and $\forall S\subset A$ with $|S|=n$, there exists a subset $T\subset S$ with
\[
|T|\ge \epsilon n
\]
and $T$ has \emph{no three collinear points}, i.e. for every affine line $\ell\subset\mathbb{R}^2$ we have $|T\cap \ell|\le 2$.

Question: Does it follow that there exists an integer $m\ge 1$ and a covering
\[
A = A_1\cup\cdots\cup A_m
\]
with each $A_i\subset\mathbb{R}^2$ having no three collinear points (equivalently: a coloring of $A$ with $m$ colors such that no line contains three points of the same color)?

Conventions/edge cases:
- A “subset of size $n$” means an $n$-element subset.
- If $A$ itself has no three collinear points then the answer is “yes” with $m=1$.
- The hypothesis forces $0<\epsilon\le 1$.
